\documentclass[aps,pre,onecolumn,superscriptaddress,final]{revtex4-2}

\usepackage{amsmath,amssymb,bm}
\usepackage[final]{graphicx}
\usepackage{physics}
\usepackage{hyperref}
\usepackage{booktabs}
\usepackage{subcaption}

% Avoid math warnings in PDF strings
\pdfstringdefDisableCommands{%
  \def\mu{mu}%
  \def\mathcal#1{#1}%
  \def\alpha{alpha}%
  \def\beta{beta}%
  \def\gamma{gamma}%
  \def\eta{eta}%
  \def\delta{delta}%
  \def\u{u}%
}

\graphicspath{{figures/}}
\setkeys{Gin}{draft=false}

\begin{document}

\title{Nonreciprocal Cancer--Immune Dynamics with Allee Effects and Immunological Control}

\author{Erick Serrato Garc\'ia}
\affiliation{Physics Department, Universidad Aut\'onoma Metropolitana Iztapalapa}

\begin{abstract}
We study a spatial cancer--immune reaction--diffusion model with a cooperative strong Allee effect, nonreciprocal couplings, and adaptive immunological control. The four strong-Allee steady-state scenarios analyzed here have linearly unstable homogeneous reference states, so the steady states organize spatial routes as unstable references rather than stable attractors. In two dimensions, initially disordered fields self-organize into mesoscopic domains whose correlation lengths grow close to a diffusion-limited $t^{1/2}$ law. The main distinction is between two reciprocity regimes. For $\mathcal{R}=0$, switching on the adaptive control $u(\mathbf{x},t)$ strongly changes the effective driving forces: the cancer potential becomes strongly negative, while the healthy-tissue and immune potentials become positive, and the correlation structure is correspondingly reorganized. For $\mathcal{R}=1$, the controlled and uncontrolled trajectories are much closer because the reciprocal-completion terms already impose a strong morphological bias; in this regime $u$ acts as a secondary correction rather than the dominant organizer. We connect these correlation signatures with a thermodynamic diagnostic based on the positive diffusive entropy-production density $\sigma^+=\sum_a D_a|\nabla\phi_a|^2$ and its integral $\Sigma_{\mathrm{diss}}(t)$. This observable quantifies gradient relaxation and the diffusive cost of interfaces, while species-resolved contributions identify whether cancer, healthy tissue, or immune structure dominates the organization. Effective potentials and their gradients provide complementary observables for comparing the scenarios. In the reciprocal completion of the dynamics, equivalent to a free-energy-minimizing gradient flow, the Hill control has no comparable effect on diffusive entropy production; its thermodynamic signature is therefore specific to the nonreciprocal regime. The simulations are reported in nondimensional model units on a domain $[0,8]^2$; physical lengths must be assigned by calibration to a histological, cellular, or tissue scale. Together, the correlation lengths, effective potentials, and diffusive entropy-production measures show that nonreciprocity and adaptive control select spatial scales and coupling structures rather than producing relaxation to a homogeneous equilibrium.
\end{abstract}

\maketitle

\section{Introduction}

\subsection{Cancer as a nonequilibrium system}
Tumor growth is a driven, open system with matter and energy fluxes. Nonequilibrium physics provides tools to interpret multistability, thresholds, and pattern formation in tumor--immune competition \cite{murray2002}.

\subsection{Limitations of reciprocal population models}
Symmetric Lotka--Volterra or logistic models miss thresholds, hysteresis, and asymmetric suppression/activation between cancer, healthy tissue, and immune cells. These limitations motivate adding Allee effects and nonreciprocal couplings.

\subsection{Allee effects and critical thresholds}
An Allee term introduces a viability threshold for cancer: below it, extinction; above it, proliferation. Weak Allee yields a soft slowdown; strong Allee creates a sharp barrier. Autocrine signalling can induce Allee-like behaviour in tumors \cite{gerlee2022}; phenotypic switching can reinforce extinction for subcritical populations \cite{bottger2015}.

\subsection{Nonreciprocal interactions in cancer--immune dynamics}
Suppression of immunity by tumor cells and cytotoxic action of immune cells are not symmetric. Including nonreciprocity alters nullclines, fixed points, and transient routes \cite{delitala2007,gatenby2004}. We treat nonreciprocity as asymmetry in interaction coefficients and as an external control field.

\subsection{Summary of contributions}
\begin{itemize}
    \item Formulate ODE/PDE models with weak/strong Allee effects, nonreciprocity, and an immunological control field $u(x,t)$.
    \item Characterize fixed points and linear stability; identify Allee-driven bifurcations in mean field.
    \item Extend to reaction--diffusion dynamics; report spatial patterns, correlation lengths, and spectral content.
    \item Evaluate control strategies (binary vs adaptive) and their cost--efficiency trade-offs across weak/strong Allee regimes.
    \item Provide figure set: Fig.~1 conceptual scheme; Fig.~2 phase planes/nullclines; Fig.~3 bifurcations; Fig.~4--5 spatial patterns and spectra; Fig.~6 control field; Table~I parameters; Table~II cost vs suppression.
\end{itemize}

% Figure 1 placeholder: conceptual scheme of cancer--immune interactions and control.
\begin{figure}[htbp]
    \centering
    \includegraphics[draft=false,width=0.55\linewidth]{steady.png}
    \caption{Conceptual phase portrait (weak Allee, $\mu=1$) illustrating nullclines and an unstable fixed point as threshold between extinction and proliferation.}
\end{figure}


\section{Model}

\subsection{Population variables}
\begin{itemize}
    \item $c(t)$ or $c(\mathbf{x},t)$: cancer cell density.
    \item $s(t)$ or $s(\mathbf{x},t)$: healthy tissue density.
    \item $i(t)$ or $i(\mathbf{x},t)$: immune effector density.
\end{itemize}

\subsection{Baseline dynamics (ODE)}
We start from logistic-like growth for each population and quadratic interaction terms capturing cytotoxicity, tissue damage, and immune cooperation \cite{delitala2007}:
\begin{align}
\dot{c} &= r_c\, c\, f_{\text{Allee}}(c) - \alpha\, c\, s - \beta\, c\, i - \tfrac{\mu}{2}(\eta\, i^2 + \gamma\, s^2), \\
\dot{s} &= r_s\, s \left(1-\tfrac{s}{k_s}\right) - \gamma\, c\, s + \delta\, s\, i^2 - \tfrac{\mu}{2}\alpha\, c^2, \\
\dot{i} &= r_i\, i \left(1-\tfrac{i}{k_i}\right) - \eta\, c\, i + \delta\, s^2 i - \tfrac{\mu}{2}\beta\, c^2 .
\end{align}

\subsection{Allee effect formulation}
\begin{itemize}
    \item \textbf{Weak Allee:} $f_{\text{Allee}}(c)=\left(\tfrac{c}{a}-1\right)\left(1-\tfrac{c}{k_c}\right)$ (soft threshold) \cite{murray2002}.
    \item \textbf{Strong Allee:} $f_{\text{Allee}}(c)=\left(\tfrac{c}{a}-1\right)\left(\tfrac{c-a}{k_c-a}\right)$ (hard threshold) \cite{gerlee2022}.
\end{itemize}

\subsection{Nonreciprocal interactions}
Suppression of immunity by cancer ($-\eta\,c\,i$) differs from immune cytotoxicity ($-\beta\,c\,i^2$); damage to healthy tissue by cancer ($-\gamma\,c\,s$) differs from competition against cancer ($-\alpha\,c\,s^2$). These asymmetries break reciprocity and reshape nullclines and bifurcations.

\subsection{Immunological control field \texorpdfstring{$u(\mathbf{x},t)$}{u(x,t)}}
\begin{itemize}
    \item Binary/parametric control: $\mu\in\{0,1\}$ toggles variational terms.
    \item Adaptive field: $u(\mathbf{x},t)=u_0\, \dfrac{c}{i+\varepsilon}$ with cap $u_{\max}$ to prioritize high-$c$, low-$i$ regions.
    \item Coupling: control can enter as enhanced cytotoxicity ($\beta\to \beta+u$) or as an external field in the $i$-equation.
\end{itemize}

\subsection{Spatial extension (PDE)}
Adding diffusion yields
\begin{align}
\partial_t c &= D_c \nabla^2 c + r_c\, c\, f_{\text{Allee}}(c) - c(\alpha s^2 + \beta i^2) - \mu c(\gamma s^2 + \eta i^2), \\
\partial_t s &= D_s \nabla^2 s + r_s s(1-s) - \gamma c^2 s + \delta i^2 s - \tfrac{\mu}{2}\alpha s c^2, \\
\partial_t i &= D_i \nabla^2 i + r_i i(1-i) + \delta s^2 i - \eta c^2 i - \tfrac{\mu}{2}\beta c^2 i .
\end{align}
Diffusion acts as a mode selector; with $D=\mathrm{diag}(D_c,D_s,D_i)$ the linear spectrum is $\lambda(k)=\mathrm{eig}(J-Dk^2)$.

\begin{table*}[!htbp]
\centering
\begin{tabular}{lll}
\toprule
Symbol & Description & Example value \\
\midrule
$r_c,r_s,r_i$ & intrinsic growth rates (cancer, healthy, immune) & $5.84,\ 13.12,\ 10.92$ \\
$k_c,k_s,k_i$ & carrying capacities & $1,\ 1,\ 1$ \\
$a$ & Allee threshold (weak/strong) & $7.22$ (weak), $0.1$ (strong scans) \\
$\alpha,\beta,\gamma,\delta,\eta$ & nonreciprocal interaction strengths & $10.22,\ 7.60,\ 0.74,\ 5.40,\ 5.08$ \\
$D_c,D_s,D_i$ & diffusion coefficients (spatial model) & $1,\ 1,\ 1$ \\
$\mu$ & control/variational switch & $\{0,1\}$ \\
$u(\mathbf{x},t)$ & adaptive immunological control field & capped by $u_{\max}$ \\
\bottomrule
\end{tabular}
\caption{Variables y par\'ametros del modelo. Ajustar valores según régimen weak/strong y escenarios de control.}
\label{tab:params}
\end{table*}


\section{Spatial dynamics}

\subsection{Reaction--diffusion extension}
Add diffusion to each population:
\begin{align}
\partial_t c &= D_c \nabla^2 c + r_c\, c\, f_{\text{Allee}}(c) - c(\alpha s^2 + \beta i^2) - \mu c(\gamma s^2 + \eta i^2), \\
\partial_t s &= D_s \nabla^2 s + r_s s(1-s) - \gamma c^2 s + \delta i^2 s - \tfrac{\mu}{2}\alpha s c^2, \\
\partial_t i &= D_i \nabla^2 i + r_i i(1-i) + \delta s^2 i - \eta c^2 i - \tfrac{\mu}{2}\beta c^2 i .
\end{align}

\subsection{Numerical methods}
Finite elements (FEniCS) in 2D; uniform mesh ($\sim$200 nodes per side), step $\Delta t \sim 5\times 10^{-3}$, final time $T\sim 2$. Random initial conditions within feasible bounds $0.8<c<1.2$, $0.2<s<0.6$, $0<i<0.2$. Snapshots y videos siguen los cuatro escenarios: weak/strong y $\mu=0/1$.

\subsection{Pattern formation}
Report snapshots of $c,s,i$ for $\mu=0$ vs $\mu=1$ and for weak vs strong Allee (Fig.~4). Extract correlation lengths and dominant wavelengths from FFT; track $\xi(t)$ from autocorrelations ($c_c,s_s,i_i$) y cruzadas ($c_s,c_i,s_i$). Compare organización y compacidad:
\begin{itemize}
  \item $\mu=0$: patrones fragmentados, dominios tumorales grandes si Allee es débil; con Allee fuerte los cúmulos son más pequeños.
  \item $\mu=1$: la estructura variacional suaviza y compacta los dominios; en Allee fuerte los tamaños quedan acotados.
  \item $\xi(t)$ crece subdifusivamente $\sim e^{\alpha} t^{0.5}$; $c$ presenta la mayor $\xi$, $s_i$ la menor.
\end{itemize}
Fig.~5 puede mostrar espectros de potencia (inicio vs fin) y curvas $\xi(t)$.

\begin{figure}[!htbp]
    \centering
    \includegraphics[draft=false,width=0.48\linewidth]{fields_mu_0.png}
    \includegraphics[draft=false,width=0.48\linewidth]{fields_mu_1.png}\\
    \includegraphics[draft=false,width=0.48\linewidth]{fields_strong_mu_0.png}
    \includegraphics[draft=false,width=0.48\linewidth]{fields_strong_mu_1.png}
    \caption{Patrones espaciales de $c,s,i$ en $T\approx1.9$: fila superior (weak Allee) $\mu=0$ y $\mu=1$; fila inferior (strong Allee) $\mu=0$ y $\mu=1$.}
\end{figure}

\begin{figure}[!htbp]
    \centering
    \includegraphics[draft=false,width=0.48\linewidth]{power_spectrum_c_0.000.png}
    \includegraphics[draft=false,width=0.48\linewidth]{power_spectrum_c_1.000.png}\\
    \includegraphics[draft=false,width=0.48\linewidth]{correlation_mean_c_c.png}
    \includegraphics[draft=false,width=0.48\linewidth]{correlation_mean_s_s.png}
    \caption{Arriba: espectro de potencia del campo $c$ al inicio y al final ($\mu=1$, weak Allee). Abajo: evolución de la longitud de correlación $\xi(t)$ para $c_c$ y $s_s$.}
\end{figure}

\begin{figure}[!htbp]
    \centering
    \includegraphics[draft=false,width=0.48\linewidth]{correlation_mean_i_i.png}
    \includegraphics[draft=false,width=0.48\linewidth]{correlation_mean_c_i.png}\\
    \includegraphics[draft=false,width=0.48\linewidth]{correlation_strong_mean_c_c.png}
    \includegraphics[draft=false,width=0.48\linewidth]{correlation_strong_mean_c_i.png}
    \caption{Correlaciones adicionales: izquierda arriba $\xi(t)$ para $i_i$, derecha arriba para $c_i$ (weak Allee); abajo correlaciones análogas en régimen Allee fuerte.}
\end{figure}

\clearpage

% Fig. 4: spatial patterns (weak/strong, mu=0/1). Fig. 5: temporal evolution of densities or spectra.


\section{Immunological control}

\subsection{Control formulation}
Represent control via $\mu$ (variational switch) and an adaptive field
\[
u(\mathbf{x},t)=\min\!\left(u_{\max},\,u_0\, \frac{c}{i+\varepsilon}\right)
\]
to prioritize high-$c$, low-$i$ regions. Coupling can enhance $\beta$ (cytotoxicity) or enter as a source in the immune equation.

\subsection{Control strategies}
\begin{itemize}
    \item Constant vs adaptive dosing.
    \item Spatial targeting: stronger control where $c$ is high and $i$ is low.
    \item Timing protocols: pulsed vs continuous application.
\end{itemize}

\subsection{Efficiency and cost}
Define a cost functional (drug load, $\int u\, d\mathbf{x}\,dt$) vs benefit (peak reduction, time to extinction, total tumor burden $\int c\,dt$, correlation length of $c$). Compare $\mu=0,1$ and adaptive $u$; report Pareto curves cost vs suppression. Qualitatively: $\mu=1$ reduces fragmentation and peak burden; adaptive $u$ focuses resources on high-$c$, low-$i$ zones, improving containment for similar cost.

\setcounter{table}{1} % ensure Table II numbering under revtex onecolumn
\begin{table*}[!htbp]
\centering
\begin{tabular}{lccc}
\toprule
Protocol & Peak $c$ reduction & Time to extinction & Cost (relative) \\
\midrule
$\mu=0$ (no control) & baseline & -- & 0 \\
$\mu=1$ (binary) & $\downarrow$ moderate & shorter vs baseline & medium \\
Adaptive $u$ (weak Allee) & $\downarrow$ strong & shortest & medium--high \\
Adaptive $u$ (strong Allee) & $\downarrow$ strong & shortest & medium--high \\
\bottomrule
\end{tabular}
\caption{Table II (qualitative): control efficiency vs cost. Replace with measured metrics (peak, extinction time, integrated $u$) when available.}
\label{tab:control-cost}
\end{table*}

\begin{figure}[!htbp]
    \centering
\includegraphics[draft=false,width=0.6\linewidth]{estabilidad_lineal_strong_mu1.png}
    \caption{Ejemplo de perfil espectral bajo control ($\mu=1$) para Allee fuerte: se atenúan modos de alta frecuencia, reduciendo fragmentación.}
\end{figure}

% Fig. 6: control field $u(x,t)$ applied on $c$-rich, $i$-poor regions.

% Fig. 6: control field $u(x,t)$. Tabla II: coste vs supresión.


\section{Thermodynamic Diagnostic}\label{sec:thermodynamics}

The transport form introduced in Sec.~\ref{sec:spatial} provides a direct nonequilibrium diagnostic for the spatial simulations. The diffusive currents $\mathbf{J}_a=-D_a\nabla\phi_a$ measure how gradients in $c$, $s$, and $i$ are relaxed by transport, while the reaction terms organize the local growth, competition, immune feedback, and control response. We therefore use the gradient content of the fields as a positive diffusive entropy-production density in model units.

In nonequilibrium thermodynamics, local dissipative densities are written as sums of flux--force products. For diffusion at uniform temperature, a classical contribution has the structure
\begin{equation}
    \sigma_S
    =-\frac{1}{T}\sum_a \mathbf{J}_a\cdot\nabla\mu_a+\cdots ,
    \label{eq:entropy-flux-force}
\end{equation}
where $\mu_a$ is a thermodynamic chemical potential. If the constitutive law is $\mathbf{J}_a=-D_a\nabla\mu_a$, then each diffusive contribution is nonnegative:
\begin{equation}
    -\mathbf{J}_a\cdot\nabla\mu_a
    =D_a\left|\nabla\mu_a\right|^2\geq 0 .
    \label{eq:positive-mu-diss}
\end{equation}

The visualization and time-series post-processing used here follow the effective convention implemented in the nonequilibrium scripts: the population fields themselves are used as the transport variables driving Fickian diffusion, i.e. $\mu_a\equiv\phi_a$ at the level of the plotted diffusive fluxes. Thus the local positive density is
\begin{equation}
    \sigma_{\mathrm{diss},a}(\mathbf{x},t)
    :=D_a\left|\nabla\phi_a(\mathbf{x},t)\right|^2,
    \label{eq:sigma-diss-field}
\end{equation}
Here $\sigma^+$ is a \emph{metric of diffusive gradients in model units}, constructed by analogy with the De Groot--Mazur entropy-production structure, but \emph{without formal identity} to the Onsager pair $-\mathbf{J}_a\cdot\nabla\mu_a$ when $\mu_a$ denotes a true thermodynamic chemical potential. The convention $\mu_a\equiv\phi_a$ is operational for visualization and cross-scenario comparison, not a strict thermodynamic claim.
The total density plotted in the spatial $\sigma^+$ maps is
\begin{equation}
    \sigma^+(\mathbf{x},t)
    \equiv\sigma_{\mathrm{diss,tot}}(\mathbf{x},t)
    =D_c|\nabla c|^2+D_s|\nabla s|^2+D_i|\nabla i|^2 .
    \label{eq:sigma-plus-density}
\end{equation}
The corresponding integrated time series is
\begin{equation}
    \Sigma_{\mathrm{diss}}(t)
    =\int_\Omega \sigma^+(\mathbf{x},t)\,dA
    =\int_{\Omega}\sum_a D_a
    \left|\nabla \phi_a(\mathbf{x},t)\right|^2\,dA .
    \label{eq:sigma-diss-integral}
\end{equation}
Introducing a constant effective temperature $T_{\mathrm{eff}}$ would only rescale this entropy-production measure by $1/T_{\mathrm{eff}}$; it would not change the spatial patterns or the relative temporal comparisons. We therefore report $\Sigma_{\mathrm{diss}}$ directly as a diffusive organization metric: large values indicate sharp interfaces and strong transport demand, whereas decay indicates smoothing and coarsening.

\begin{figure*}[!htbp]
    \centering
    \begin{subfigure}[t]{0.48\textwidth}
        \centering
        \includegraphics[draft=false,width=\linewidth]{noneq_sigma_diss_comparison.png}
        \caption{$\Sigma_{\mathrm{diss}}(t)$}
        \label{fig:noneq_sigma_diss}
    \end{subfigure}\hfill
    \begin{subfigure}[t]{0.48\textwidth}
        \centering
        \includegraphics[draft=false,width=\linewidth]{noneq_sigma_c_comparison.png}
        \caption{$\Sigma_c=\int D_c|\nabla c|^2dA$}
        \label{fig:noneq_sigma_c}
    \end{subfigure}

    \vspace{0.6em}

    \begin{subfigure}[t]{0.48\textwidth}
        \centering
        \includegraphics[draft=false,width=\linewidth]{noneq_sigma_s_comparison.png}
        \caption{$\Sigma_s=\int D_s|\nabla s|^2dA$}
        \label{fig:noneq_sigma_s}
    \end{subfigure}\hfill
    \begin{subfigure}[t]{0.48\textwidth}
        \centering
        \includegraphics[draft=false,width=\linewidth]{noneq_sigma_i_comparison.png}
        \caption{$\Sigma_i=\int D_i|\nabla i|^2dA$}
        \label{fig:noneq_sigma_i}
    \end{subfigure}
    \caption{Diffusive entropy-production diagnostics for the strong-Allee scenarios in the window $0\leq t\leq 1$, sampled with $\Delta t=10^{-3}$. The total integrated entropy production $\Sigma_{\mathrm{diss}}(t)$ is shown together with the cancer, healthy-tissue, and immune contributions. Larger values correspond to sharper interfaces and stronger spatial gradients; decay indicates diffusive smoothing and coarsening. The uncontrolled nonreciprocal baseline retains the largest late-time gradient content, whereas Hill control strongly suppresses the cancer contribution and the reciprocal-completion curves remain close with or without control, indicating that this setting has already selected a similar spatial route.}
    \label{fig:noneq_sigma_panel}
\end{figure*}

The species-resolved panels in Fig.~\ref{fig:noneq_sigma_panel} give the temporal chronology of gradient relaxation. Early times are dominated by sharp interfaces generated from the initial heterogeneous texture. As coarsening proceeds, each compartment relaxes on its own scale, revealing whether the spatial organization is controlled by tumor fronts, healthy-tissue recovery, or immune localization.

We also compute effective potentials from the variational post-processing functional,
\begin{equation}
    \mu_a^{\mathrm{eff}}(\mathbf{x},t)
    \sim \frac{\delta F}{\delta\phi_a},
    \label{eq:effective-potentials}
\end{equation}
and use their spatial averages to compare the mean direction and intensity of the population-level driving forces. Because the reaction Jacobian can be decomposed as
\begin{equation}
    A_{ab}:=\frac{\partial R_a}{\partial\phi_b},
    \qquad
    A=S+N,
    \label{eq:jacobian-symmetric-antisymmetric}
\end{equation}
with a generally nonzero antisymmetric part $N$, the homogeneous reaction subsystem contains a rotational or nonreciprocal component. We therefore use the effective potentials as comparative observables for the modeled population dynamics.

This decomposition also clarifies the reciprocal reference case. If the reaction dynamics can be completed as a gradient flow,
\begin{equation}
    R_a(\phi)=-M_a\,\frac{\delta F}{\delta\phi_a},
    \qquad M_a>0,
    \label{eq:reciprocal-gradient-flow}
\end{equation}
then the antisymmetric part vanishes ($N=0$) and the free-energy functional decreases monotonically,
\begin{equation}
    \frac{dF}{dt}
    =-\sum_a\int_\Omega M_a
    \left|\frac{\delta F}{\delta\phi_a}\right|^2 dA
    -\sum_a\int_\Omega D_a
    \left|\nabla\frac{\delta F}{\delta\phi_a}\right|^2 dA
    \leq 0 .
    \label{eq:free-energy-descent}
\end{equation}
The nonreciprocal cases analyzed below should be read against this gradient-flow baseline: they need not minimize a single global $F$, and the Hill control can alter the spatial redistribution of gradients rather than simply accelerating descent along a free-energy landscape.

\begin{figure*}[!htbp]
    \centering
    \begin{subfigure}[t]{0.48\textwidth}
        \centering
        \includegraphics[draft=false,width=\linewidth]{noneq_mu_c_effective_comparison.png}
        \caption{$\langle\mu_c^{\mathrm{eff}}\rangle$}
        \label{fig:noneq_mu_c}
    \end{subfigure}\hfill
    \begin{subfigure}[t]{0.48\textwidth}
        \centering
        \includegraphics[draft=false,width=\linewidth]{noneq_mu_s_effective_comparison.png}
        \caption{$\langle\mu_s^{\mathrm{eff}}\rangle$}
        \label{fig:noneq_mu_s}
    \end{subfigure}

    \vspace{0.6em}

    \begin{subfigure}[t]{0.48\textwidth}
        \centering
        \includegraphics[draft=false,width=\linewidth]{noneq_mu_i_effective_comparison.png}
        \caption{$\langle\mu_i^{\mathrm{eff}}\rangle$}
        \label{fig:noneq_mu_i}
    \end{subfigure}
    \caption{Spatial averages of the effective cancer, healthy-tissue, and immune potentials for the four strong-Allee scenarios. The panel highlights how control and reciprocity structure modify the mean driving forces: in the nonreciprocal baseline, Hill control shifts the cancer potential from weakly favorable to strongly suppressive and raises the healthy-tissue and immune potentials; in the reciprocal-completion setting, controlled and uncontrolled trajectories nearly overlap, showing that the additional coupling structure already imposes the dominant thermodynamic bias.}
    \label{fig:noneq_mu_panel}
\end{figure*}

The effective potentials make the contrast between control regimes explicit. For $\mathcal{R}=0$, activating the adaptive control changes the final averaged potentials from a weak cancer-favorable balance, approximately $\langle\mu_c^{\mathrm{eff}}\rangle\simeq 6.2\times10^{-2}$, $\langle\mu_s^{\mathrm{eff}}\rangle\simeq -1.7\times10^{-2}$, and $\langle\mu_i^{\mathrm{eff}}\rangle\simeq -1.4\times10^{-1}$, to a strongly cancer-suppressing and healthy--immune-favoring state, approximately $\langle\mu_c^{\mathrm{eff}}\rangle\simeq -19.7$, $\langle\mu_s^{\mathrm{eff}}\rangle\simeq 6.1$, and $\langle\mu_i^{\mathrm{eff}}\rangle\simeq 6.2$. This explains why the $\mathcal{R}=0$ correlations change substantially when $u$ is switched on: the control reverses the effective population driving forces and changes which fields organize the interfaces.

For $\mathcal{R}=1$, the uncontrolled and controlled trajectories are much closer. Both end with strongly negative cancer potential and positive healthy/immune potentials, approximately $\langle\mu_c^{\mathrm{eff}}\rangle\simeq -35$, $\langle\mu_s^{\mathrm{eff}}\rangle\simeq 7.47$, and $\langle\mu_i^{\mathrm{eff}}\rangle\simeq 7.17$. This happens because the $\mathcal{R}=1$ scenarios already include the reciprocal-completion terms and use the stronger cancer--immune coupling set ($\beta=5$ instead of $\beta=3$, with $r_d=10$ instead of $14$). In this regime, reciprocity completion has already selected a cancer-suppressed, healthy--immune organized route, so adaptive control produces a smaller additional displacement in the effective-potential and correlation curves.

Thus, the large separation between Hill-on and Hill-off curves for $\mathcal{R}=0$ arises because the control supplies the dominant geometric bias in a nonreciprocal background. By contrast, for $\mathcal{R}=1$, the reciprocal-completion terms already impose a cancer-suppressing, healthy--immune organized route; Hill control therefore produces only a secondary correction in both effective potentials and diffusive entropy production.

In the reciprocal completion of the dynamics, equivalently the gradient-flow limit in Eqs.~\eqref{eq:reciprocal-gradient-flow}--\eqref{eq:free-energy-descent}, the Hill control has no comparable effect on the diffusive entropy production. Its thermodynamic signature is therefore tied to the nonreciprocal regime, where the control changes the spatial redistribution of gradients rather than simply relaxing along a free-energy descent.

For completeness, we also evaluate the gradient content of these effective potentials:
\begin{equation}
    \Sigma_{\mu}(t)=
    \int_{\Omega}\sum_a D_a
    \left|\nabla \mu_a^{\mathrm{eff}}(\mathbf{x},t)\right|^2\,dA .
    \label{eq:sigma-mu}
\end{equation}
In general $\Sigma_\mu\neq\Sigma_{\mathrm{diss}}$, unless the effective potentials are linearly related to the fields in the relevant range. The primary observable in this paper remains $\sigma^+=\sigma_{\mathrm{diss,tot}}$, while $\Sigma_\mu$ is used as a complementary measure of heterogeneity in the effective driving forces.


\section{Discussion}\label{sec:discussion}

\subsection{Physical interpretation of results}
Our results highlight an intrinsically steady-state-centered organization of cancer--immune dynamics: in the strong-Allee scenarios considered here, relevant homogeneous states act as unstable references rather than attractors. This is reflected by Re\,$\lambda_{\max}>0$ in the local analysis; in space, the same instability manifests as pattern selection and coarsening rather than convergence to a homogeneous equilibrium. These steady states organize spatial routes, while nonreciprocity tilts the phase flow and biases the geometry of cancer--healthy--immune tracking.

Two complementary observables make this organization explicit. First, field snapshots in temporal panels (Fig.~\ref{fig:fields_panel_strong_mu1_u}) show a rapid transition from disordered initial conditions to mesoscopic domains whose morphology depends on the control protocol. Second, correlation length grids $\xi(t)$ (Fig.~\ref{fig:corr_grid_auto_cross}) quantify the \emph{emerging length scale} of these domains, separating single-field coherence from inter-field co-localization. Across scenarios, $\xi(t)$ grows subdiffusively, consistent with diffusion-limited coarsening, and is well described by $\xi(t)\sim e^{\alpha}t^{1/2}$ over the simulated horizon. The thermodynamic diagnostic in Sec.~\ref{sec:thermodynamics} complements this picture through the diffusive entropy production: $\Sigma_{\mathrm{diss}}(t)$ is largest when interfaces are sharp and decreases as the fields smooth and coarsen. Its species-resolved components identify which population field controls the remaining spatial organization.

\subsection{Role of \texorpdfstring{$\mathcal{R}$}{R} (variational and reciprocity structure)}
The parameter $\mathcal{R}$ acts as a structural switch that reorganizes the spatial spectrum: for $\mathcal{R}=1$ patterns are typically less fragmented, with smoother interfaces and fewer high-frequency features, consistent with attenuation of short-wavelength modes. This qualitative observation is supported by correlation length grids, where $\mathcal{R}$ shifts $\xi(t)$ systematically in the strong Allee scenarios (Fig.~\ref{fig:corr_grid_auto_cross}). The effect is strongest in immune-centered observables: $i$--$i$, $c$--$i$, and $s$--$i$ increase their prefactors when $\mathcal{R}=1$, indicating more coherent immune domains and stronger spatial tracking of tumor and tissue structure. Importantly, $\mathcal{R}$ does not stabilize dynamics in the tested parameter set; instead it functions as a \emph{morphological selector} that changes the geometry of transient trajectories and the characteristic scale of domains.

\subsection{Adaptive immunological control \texorpdfstring{$u(\mathbf{x},t)$}{u(x,t)} as geometric coupling}
Adaptive control introduces a spatially directed action that is stronger in regions with high $c$ and low $i$. Field panels (Fig.~\ref{fig:fields_panel_strong_mu1_u}) show that activating/deactivating $u$ changes the route by which domains form and fuse. Quantitatively, $c$--$c$ tracks coherence of tumor domains, while $c$--$i$ tracks cancer--immune co-localization, i.e., the degree to which immunity follows tumor structure. The two-column organization in Fig.~\ref{fig:corr_grid_auto_cross} separates these effects visually: autocorrelations measure domain compactness within each population, whereas cross-correlations measure the geometry of interaction between populations. In the reciprocal completion of the dynamics, equivalent to free-energy minimization, the same control has no comparable impact on diffusive entropy production; its thermodynamic signature is therefore a genuinely nonreciprocal effect. This provides a concrete and measurable bridge between mechanistic control rules and emerging spatial organization.

\subsection{Comparison with previous models}
Relative to reciprocal Lotka--Volterra models and classical immune--tumor ODE/PDE models without Allee terms, the key qualitative difference is that a \emph{cooperative Allee effect} introduces a hard/smooth threshold separating collapse from proliferation, while nonreciprocity breaks the symmetry of interaction terms and biases transient routes. In spatial configurations, this combination shifts the pattern formation landscape from purely diffusion-driven instabilities to \emph{threshold-mediated} domain formation, where Allee type establishes the ease of domain expansion (weak) or imposes a critical barrier limiting domain sizes (strong). Control protocols then act mainly by selecting spatial scale and coupling structure, rather than creating stable homogeneous coexistence.

\subsection{Implications and limitations}
The present scenario set supports a robust qualitative conclusion: within explored parameters, the system operates near unstable steady states and the biologically relevant outcome is controlled \emph{containment vs escape} rather than asymptotic stability. Stability may still emerge in extended sweeps (e.g., reduced $r_c$ or improved immune recruitment via larger $r_d,\delta$), but spatial results already show that control and Allee strength strongly constrain morphology and characteristic length scales.

Limitations follow naturally. (i) The spatial horizon is short ($T=1$ in model units), so reported scaling and morphologies should be interpreted as transient-but-organized regimes. (ii) We report one realization per scenario (nb=1), so ensemble variability is not yet quantified; the correlation length framework provides a direct route to such ensemble statistics in future work. (iii) Adaptive control may introduce numerical stiffness in some runs (diagnostic triggers), suggesting that solver tolerances and control smoothness must be co-designed to ensure robust long-time integration.

\section{Conclusions}

We have analyzed a spatial cancer--healthy--immune reaction--diffusion model in which cooperative strong-Allee growth, nonreciprocal interactions, and Hill-type immunological control act together to select mesoscopic organization. The homogeneous reference states considered here are linearly unstable, so they do not serve as attracting equilibria. Instead, they organize spatial routes: initially heterogeneous fields coarsen into domains whose correlation lengths grow close to a diffusion-limited $t^{1/2}$ law over the simulated window.

The reciprocity parameter $\mathcal{R}$ is the main structural selector. In the nonreciprocal baseline ($\mathcal{R}=0$), activating Hill control supplies the dominant geometric bias: it suppresses tumor-domain coherence, reorganizes cancer--immune co-localization, and produces large changes in the effective cancer, healthy-tissue, and immune potentials. In the reciprocal-completion setting ($\mathcal{R}=1$), the coupling structure already selects a cancer-suppressing, healthy--immune organized route. Consequently, turning Hill control on produces only a secondary correction in both the effective potentials and the diffusive entropy production.

The thermodynamic diagnostic $\sigma^+=\sum_a D_a|\nabla\phi_a|^2$ and its integral $\Sigma_{\mathrm{diss}}(t)$ quantify this distinction at the level of gradients and interfaces. The control has a strong entropy-production signature only when it acts on a nonreciprocal background; in the reciprocal-completion limit, equivalent to free-energy minimization, the same control loses comparable thermodynamic relevance. Thus the principal outcome is not relaxation to a homogeneous equilibrium, but selection of spatial scales, routes, and coupling structures.

These results provide a compact framework for comparing intervention strategies through mesoscopic observables: correlation lengths, effective potentials, and diffusive entropy-production measures. Future work should extend the parameter sweeps, optimize the Hill control in space and time, test ensemble variability across initial conditions, and calibrate the nondimensional length scale against histological or cellular data.


\bibliography{bibliography}

\end{document}

 